\section{A concrete Lean route, and its present status}
\label{sec:lean-route}

The relevant target judgment in the inspected repository is
\begin{center}
\texttt{Prf PRAAx \(\Gamma\) \(\varphi\)},
\end{center}
where \texttt{PTm n} is the syntax of primitive-recursive terms and
\texttt{Fm PTm n} is the syntax of formulas in a common finite scope.
This target must be kept distinct from a theorem about Lean's natural
numbers, rationals, reals, or matrices.  The existing bridge translates
ordinary ConciseIPRA proofs to the Foundation presentation of ordinary
iPRA; the bridge does not convert arbitrary host-level mathematics into
object-theory proofs.

\paragraph{Verification status.}
All repository APIs discussed here were read at revision
\path{3ee850b4c2d650cf01263823fd5415963747c5d0}.
That revision requests Lean \texttt{v4.32.0}.
No Lean executable was available in the preparation environment.
The new Lean source is therefore an \emph{unelaborated formalization
outline}, including a complete proposed proof of the finite cancellation
lemma below.  No new kernel-checked Lean theorem is claimed.
The three runtime norm targets are definitions of remaining obligations,
not admitted theorems.  The absence of an admitted-proof command in those
files is not evidence that the obligations have been discharged.

\subsection{An object-theory lemma with an explicit certificate}

The repository represents an integer term by two natural terms
\(u^+,u^-\), with intended value \(u^+-u^-\).
Its integer comparison \texttt{zle u v} expands to the natural comparison
\(u^++v^-\leq v^++u^-\).
The module \texttt{Core/IntOrder.lean} associates the natural slack
\[
 s(u,v)=(v^++u^-)\mathbin{\dot-}(u^++v^-)
\]
with a proof of this comparison.  Here \(\dot-\) is truncated subtraction.
Thus a signed comparison is usable through an equality with an explicit
nonnegative witness, without deciding the sign of an arbitrary real.

Suppose signed numerators \(A,B,C\) and natural terms \(w,P,e_L,e_R\)
satisfy
\[
 A\leq C+e_L,\qquad B\leq C+e_R,\qquad
 2wC+P\leq w(A+B).
\]
Let their respective slacks be \(s_1,s_2,s_3\).  The whole argument is
the polynomial identity
\[
 w(e_L+e_R)=P+s_3+w(s_1+s_2).
\]
Its right-hand remainder is a natural term, so it proves
\(P\leq w(e_L+e_R)\) in iPRA.
The file \texttt{MidpointCancellation.lean} implements this certificate
using the existing \texttt{zle\_toZ}, \texttt{le\_of\_eq\_add}, and
\texttt{deriveNat} interfaces.  The latter interface reflects a proved
identity in the ring of provable term-equivalence classes back to an
object equation.  It is not reflection from arbitrary truth in the
standard model.

For the cloud theorem, use
\[
 A=2F_z(x),\quad B=2F_z(y),\quad
 C=2F_z((x+y)/2),\quad w=1,\quad P=\|x-y\|_2^2.
\]
The parallelogram identity supplies the third premise.  The two repaired
midpoints supply the first two with
\[
 e_L=L(2\eta_S+h)+2\epsilon_S,\qquad
 e_R=L(2\eta_T+h)+2\epsilon_T.
\]
The resulting object inequality has the exact constant
\[
 \|x-y\|_2^2\leq
 2L(\eta_S+\eta_T+h)+2(\epsilon_S+\epsilon_T).
\]
All rational quantities must first be placed on a common positive
denominator.  The source gives the integral cancellation step; rational
scaling and the geometric hypotheses remain explicit adapter obligations.
The identity itself was independently expanded to zero with exact
symbolic arithmetic.  That check does not constitute Lean elaboration.

\subsection{The finite-cloud implementation}

The dimension, cloud lengths, point coordinates, precision indices, and
error bounds are numerical inputs.  Fix the finite-list and rational
codecs once.  A vector is a numerical list of signed rational coordinate
codes; a cloud is a nonempty numerical list of vectors.  The assertions
of a box bound, sup-norm proximity, Hausdorff proximity, approximate
midpoint closure, and approximate minimality are bounded formulas over
these lists.  Their Boolean translations, rather than the displayed
quantified formulas themselves, are used as induction invariants.

The primitive-recursive selector chooses the least cost-minimizing
index.  The repository already supplies
\texttt{ApproximateWitness.argmin}, its range and attainment theorem,
and \texttt{argmin\_minimal}.  All finitely many costs can be converted
to natural numerators on a common positive denominator.  The selector
therefore compares exact integers; no real comparison is hidden in it.
The actual nearest-neighbor and Hausdorff scans in
\texttt{FiniteHausdorff.lean} can similarly be reused after proving the
chosen vector metric and scaling adapter.

There are four new geometric verification loops.  First, accumulate the
cost over coordinates.  The invariant says that the current accumulator
equals the sum of the first \(i\) squared differences.  Second, prove
the Lipschitz estimate by accumulating the first \(i\) coordinate
estimates, with bound \(4iBh\).  Third, accumulate the coordinatewise
parallelogram identities.  Fourth, apply the bounded midpoint and
Hausdorff witnesses and the argmin theorem to obtain the two comparisons
used above.  All loop lengths are displayed numerical bounds; even a
cubic exhaustive midpoint scan is primitive recursive.

The regular-name theorem then contains only finite comparisons at two
independent precision indices and the explicit precision schedule.
The usual interpretation as a completed real vector belongs to the
external semantic explanation.  No completion, compactness, or infinite
choice axiom is added to the arithmetic theory.

\subsection{The runtime norm targets}

The source \texttt{GaussianPolynomialMonicNorm.lean} already computes
an actual polynomial-valued multiplication matrix and its determinant.
Its inputs \(F,G,m,a,c\) are numerical: \(m\) is not a host-level matrix
dimension.  The stored norm has coarse degree cap
\[
 m\bigl(a+(m-1)c\bigr).
\]
The inspected implementation proves coefficient and evaluation formulas
and binds the actual returned matrix to the source curve equations.
It explicitly leaves norm nonvanishing, sharp degree, the runtime
matrix-kernel implication, fiber multiplicity, and the general coded-field
adapter as further obligations.

The coefficient convention matters.  The array \(F\) stores only the
lower coefficients of a monic polynomial in \(Y\), in descending
\(Y\)-order.  Each row is itself a descending coefficient list in
\(X\).  Row \(i<m\), slot \(j\leq c\), therefore corresponds to
\(X^{c-j}Y^{m-1-i}\).  The shape predicate fixes each used row's width;
it does not require raw rational codes to be canonical.  Weighted support
is consequently tested after coefficient normalization.

The file \texttt{RuntimeNormTargets.lean} defines the exact closed
object-theory target saying that a common zero of the original monic
polynomial and \(G\) is a zero of the actual computed norm.  This target
does not accept a determinant-kernel theorem as an input assumption.
The existing \texttt{returned\_zero\_action} gives every coordinate of
the relevant left-kernel equation.  The remaining finite determinant
argument must use the actual numerical determinant and the fact that
the first evaluation-vector coordinate is one.

For the weighted-degree target, use integer weights
\(w_X\geq1\), \(w_Y\geq0\).  The finite support requirements are
\[
 w_X(c-j)+w_Y(m-1-i)\leq mw_Y
\]
for nonzero lower coefficients of \(F\), and the analogous inequality
with cap \(D\) for \(G\).  The conclusion states that the actual norm's
leading slots above \(\lfloor mD/w_X\rfloor\) vanish.  It is a support
statement, not a claim that the already allocated buffer gets shorter.

The required induction invariant for the \(j\)-th companion iterate is
\[
 w_X\deg_X h_r+w_Yr\leq D+w_Yj,
\]
for each nonzero coefficient polynomial \(h_r(X)\) of \(Y^r\).
Use bounded coefficient convolution and the monic relation to prove one
step.  Apply quantifier-free induction to the Booleanized support
predicate over the actual column loop.  In each determinant term the
row-weight sum equals the column-weight sum because the numeric word
encodes a permutation.  The \(w_Y\)-contributions cancel; hence every
term has \(X\)-degree at most \(\lfloor mD/w_X\rfloor\).
A second bounded induction over the permutation scan proves the leading
zero coefficients of the stored determinant.

The repository's fixed-size determinant laws give useful mathematical
templates.  They are not already the numerical-dimension theorem just
described.  The runtime scan, its coverage and uniqueness, and an actual
adjacent-column alternation implementation are available and should be
the starting point.  The archive's formalization plan lists the exact
source paths and the remaining steps, including the separate
nonvanishing and fiber-multiplicity arguments.

\subsection{A finite-place certificate for nonvanishing}

The third target, \texttt{LocalNewtonNormTarget}, uses the existing
numeric translation program
\texttt{GaussianRationalTranslation.translated}.
For a stored polynomial of degree cap \(C\), its entry at \(C-i\)
is the coefficient of \(T^i\) in \(P(a+T)\).
The predicate \texttt{exactOrder} requires a specified nonzero leading
coefficient and zero coefficients at every smaller index.
All these checks are finite; no formal-series object or infinity-valued
valuation is a term of the arithmetic language.

The local certificate for
\(F=Y^m+\sum_{j<m} f_j(X)Y^j\) gives \(e\geq1\),
\(\gcd(m,e)=1\), and \(c_0\ne0\), such that \(f_0(a+T)\)
has order \(e\) and leading coefficient \(-c_0\).
For \(0<j<m\), every coefficient of \(T^i\) with
\(mi+ej\leq me\) vanishes.  This finite support condition handles a
zero coefficient polynomial without a special valuation convention.
For reduced \(G\), a second certificate identifies a nonzero coefficient
\(bT^vY^k\), \(k<m\), at minimal weight \(s=mv+ek\).
The coprimality condition makes the minimal pair unique.

The proposed endpoint asserts that the actual computed norm has exact
order \(s\) at \(a\), with leading coefficient
\[
 b^m\bigl((-1)^{m-1}c_0\bigr)^k.
\]
It also gives the independent global degree bound under the global
weighted-support hypotheses.  Both conclusions refer to the original
numeric norm program; the desired nonvanishing statement is not an
input premise.

For formalization, one need not introduce Puiseux roots.  At initial
weight the reduction relation is \(Y^m=c_0T^e\).  In multiplication
column \(j\), the initial monomial goes to row
\((k+j)\bmod m\), with coefficient
\(b c_0^{\lfloor(k+j)/m\rfloor}\) and \(T\)-exponent
\(v+e\lfloor(k+j)/m\rfloor\).
This row pattern is a permutation.  Its product has order \(mv+ek\),
and its cyclic-shift sign is \((-1)^{(m-1)k}\).
For any other contribution, the strict weight excess survives after the
row and column weights cancel.  Finite support induction, the actual
permutation codec, and coefficient convolution therefore provide a direct
path to the displayed coefficient formula.  Proving these connections in
Lean remains work identified by the target, not work reported as checked.

\subsection{Formalizing the polynomial-division nonvanishing certificate}
\label{sec:squarefree-lean}

The squarefree norm theorem admits a direct numerical formalization that
does not construct roots or formal branches.  Its input is a list of
exact coefficient codes for \(g,h,A_0,\ldots,A_{m-1}\), numerical
\(m,e\geq1\), and finite certificates
\[
 f=g^e h,\qquad \alpha g+\beta h=1,\qquad
 Ug+Vg'=1,\qquad \gcd(m,e)=1,\qquad\deg g\geq1.
\]
Coefficient equality and zero tests mean equality after normalization,
checked over an explicitly padded finite list.  The theorem to be
exported concerns the existing numeric norm program for \(Y^m-f\).
The ordinary Mathlib goals in the archive are semantic targets; they do
not themselves have the object-theory type \texttt{Prf PRAAx}.
For the existing constructor, store \(m\) descending \(Y\)-rows for
\(F\), all zero except the last row \(-f\), and store
\(A_{m-1},\ldots,A_0\) for \(G\).  Pad each row to a common
displayed \(X\)-degree cap.  These finite constructions discharge the
existing row-shape guards.

\paragraph{The producer and its bounded invariant.}
General long division over the coefficient field can be implemented in
\(\operatorname{degcap}(A)+1\) steps; an arbitrary nonzero leading coefficient is
handled by exact rational inversion.  Repeating division by \(g\) uses
another loop with the same cap, freezing at the first nonzero remainder.
For nonzero \(A\), the active-state invariant is
\[
 A=g^aB,\qquad B\ne0,\qquad
 \deg B+a\deg g=\deg A.
\]
This proves termination because \(\deg g\geq1\).
A separate zero flag handles \(A=0\); no infinity-valued order is
encoded.  All degrees are bounded scans of normalized coefficients.
The resulting factors \(A_j=g^{a_j}B_j\) determine the least weight
\(s=ma_k+ek\), using a finite scan over nonzero indices.

\paragraph{The actual determinant identity.}
In column \(c\), reduction of \(A_jY^{j+c}\) uses
\(j+c=r+m\delta\), with \(r<m\) and \(\delta\in\{0,1\}\).
Its matrix entry is \(A_jf^\delta\).  For each numerical permutation
word, summing this identity gives
\(\sum_c j_c=m\sum_c\delta_c\).
Thus its \(g\)-exponent \(E\) satisfies
\[
 mE=\sum_c(ma_{j_c}+ej_c)\geq ms.
\]
Equality forces every \(j_c=k\), so the unique minimal term is the
cyclic row permutation \(c\mapsto(k+c)\bmod m\).  Its carries sum to
\(k\), and its sign is \((-1)^{k(m-1)}\).
All other nonzero terms have \(E\geq s+1\).  The remainder polynomial
\(Q\) is therefore constructed directly by summing their signed
products with exponents \(E-s-1\).  This proves the finite certificate
\[
 N=g^s\left((-1)^{k(m-1)}B_k^m h^k+gQ\right).
\]
The existing permutation scan has horizon \(m^m\).
Prefix identities for products and sums, together with the actual
coefficient Leibniz theorem, provide the quantifier-free induction route.

\paragraph{Reducedness by explicit Bezout identities.}
To prove \(g\mid B^r\Rightarrow g\mid B\), compute
\(D=\gcd(g,B)\), write \(g=DC\), \(B=DE\), and retain an extended
gcd certificate.  It yields a Bezout identity for \(C,E\).
Substituting \(g=DC\) into \(Ug+Vg'=1\) gives
\[
 C(UD+VD')+D(VC')=1.
\]
Multiplying these two identities gives \(HC+JB=1\).
If \(C\mid B^r\), the finite identity
\[
 1-(JB)^r=HC\sum_{i<r}(JB)^i
\]
exhibits an inverse for \(C\).  Consequently \(g\mid B\).
The zero-polynomial case is immediate; all other cancellations have
explicit nonzero factors.  This is a finite polynomial-gcd proof of the
required reduced quotient property, with no factorization step.

Because \(B_k\) has nonzero remainder and \(h\) is invertible modulo
\(g\), the parenthesized factor in the norm identity has nonzero
remainder.  It follows that the norm is nonzero and its exact
\(g\)-divisibility order is \(s\).  A bounded scan of its actual output
buffer then returns a nonzero coefficient: if the scan failed, every
coefficient would be zero, contradicting the residue calculation.

For squarefree nonconstant rational \(f\), take \(g=f,e=1,h=1\).
The extended gcd of \(f,f'\) supplies the reducedness certificate.
The new Lean obligations are the general polynomial-division and gcd
loops, their invariants, the runtime carry/sign proof, and binding the
coefficient identity to the existing norm output.  Existing synthetic
division and monic reduction are useful templates; this report does not
assert that a general-divisor API is already available.
The attached source and roadmap remain uncompiled formalization outlines.
